% GENERATED FILE. Edit canonical lessons, then run npm run generate:books.
% canonical-source-hash: fnv1a64:d84d06c272012583
% canonical-lessons: LA-C70-filius, LA-C70-filia, LA-C70-avus, LA-C70-avia, LA-C70-familia

\chapter{Son, Daughter, Grandparents}
\label{ch:la-son-daughter-grandparents}
\hlchaptermodality{70}

\begin{chapteropening}
\textbf{By the end of this chapter:} \emph{I can name a son, a daughter, a grandfather, a grandmother and a family.}

Complete the last lesson of chapter 70: I can name a son, a daughter, a grandfather, a grandmother and a family.
\end{chapteropening}

\section[fīlius, fīliī]{fīlius, fīliī — a son}

\label{lesson:LA-C70-filius}



\begin{quote}
Before the new one: say the Latin for happy, then the Latin for sad.
\end{quote}

\subsection*{You'll want to know: fīlius, fīliī}
\textbf{fīlius, fīliī} — \textquotedblleft{}a son\textquotedblright{}. English \emph{filial}.

\begin{cousinweb}
\textbf{fīlius} is inherited. Old Latin had \emph{fīlios}, \textquotedblleft{}son\textquotedblright{}, and the word is usually traced to Proto-Indo-European \emph{*dheh1ylios}, \textquotedblleft{}a suckling\textquotedblright{} (EtymDB 2.1), from a root for nursing.
\end{cousinweb}

\begin{grammarlens}[title={What you've built}]
One more word for this chapter.
\end{grammarlens}

\subsection*{Your turn}
\begin{itemize}
\raggedright
  \item \emph{Say it:} \emph{fīlius}
  \item \emph{Say it:} \emph{fīlius}, once more
  \item \emph{Say it:} say \emph{fīlius}
  \item \emph{Recall:} say the Latin for happy, then the Latin for sad, then say \emph{fīlius} again
\end{itemize}

\begin{tcolorbox}[breakable,colback=teal!4,colframe=teal!35!black,title={Before you move on}]
What does fīlius, fīliī mean? (\textquotedblleft{}A son\textquotedblright{}.) Say it once more.
\end{tcolorbox}
\section[fīlia, fīliae]{fīlia, fīliae — a daughter}

\label{lesson:LA-C70-filia}



\begin{quote}
Before the new one: say the Latin for sad, then the Latin for a son.
\end{quote}

\subsection*{You'll want to know: fīlia, fīliae}
\textbf{fīlia, fīliae} — \textquotedblleft{}a daughter\textquotedblright{}.

\begin{cousinweb}
\textbf{fīlia} is the feminine of \emph{fīlius}, \textquotedblleft{}a son\textquotedblright{} (Lewis \& Short); the pair is usually traced to Proto-Indo-European \emph{*dheh1ylios}, \textquotedblleft{}a suckling\textquotedblright{} (EtymDB 2.1).
\end{cousinweb}

\begin{grammarlens}[title={What you've built}]
One more word for this chapter.
\end{grammarlens}

\subsection*{Your turn}
\begin{itemize}
\raggedright
  \item \emph{Say it:} \emph{fīlia}
  \item \emph{Say it:} \emph{fīlia}, once more
  \item \emph{Say it:} say \emph{fīlia}
  \item \emph{Recall:} say the Latin for sad, then the Latin for a son, then say \emph{fīlia} again
\end{itemize}

\begin{tcolorbox}[breakable,colback=teal!4,colframe=teal!35!black,title={Before you move on}]
What does fīlia, fīliae mean? (\textquotedblleft{}A daughter\textquotedblright{}.) Say it once more.
\end{tcolorbox}
\section[avus, avī]{avus, avī — a grandfather}

\label{lesson:LA-C70-avus}



\begin{quote}
Before the new one: say the Latin for a son, then the Latin for a daughter.
\end{quote}

\subsection*{You'll want to know: avus, avī}
\textbf{avus, avī} — \textquotedblleft{}a grandfather\textquotedblright{}.

\begin{cousinweb}
\textbf{avus} is inherited, usually traced to Proto-Indo-European \emph{*h2ewh2os}, \textquotedblleft{}grandfather, elder male relative\textquotedblright{} (EtymDB 2.1). Its family includes \emph{avunculus}, \textquotedblleft{}mother's brother\textquotedblright{}, the source of English \emph{uncle}.
\end{cousinweb}

\begin{grammarlens}[title={What you've built}]
One more word for this chapter.
\end{grammarlens}

\subsection*{Your turn}
\begin{itemize}
\raggedright
  \item \emph{Say it:} \emph{avus}
  \item \emph{Say it:} \emph{avus}, once more
  \item \emph{Say it:} say \emph{avus}
  \item \emph{Recall:} say the Latin for a son, then the Latin for a daughter, then say \emph{avus} again
\end{itemize}

\begin{tcolorbox}[breakable,colback=teal!4,colframe=teal!35!black,title={Before you move on}]
What does avus, avī mean? (\textquotedblleft{}A grandfather\textquotedblright{}.) Say it once more.
\end{tcolorbox}
\section[avia, aviae]{avia, aviae — a grandmother}

\label{lesson:LA-C70-avia}



\begin{quote}
Before the new one: say the Latin for a daughter, then the Latin for a grandfather.
\end{quote}

\subsection*{You'll want to know: avia, aviae}
\textbf{avia, aviae} — \textquotedblleft{}a grandmother\textquotedblright{}.

\begin{cousinweb}
\textbf{avia} is the feminine of \emph{avus}, \textquotedblleft{}a grandfather\textquotedblright{} (Lewis \& Short), which is usually traced to Proto-Indo-European \emph{*h2ewh2os}, \textquotedblleft{}grandfather\textquotedblright{} (EtymDB 2.1).
\end{cousinweb}

\begin{grammarlens}[title={What you've built}]
One more word for this chapter.
\end{grammarlens}

\subsection*{Your turn}
\begin{itemize}
\raggedright
  \item \emph{Say it:} \emph{avia}
  \item \emph{Say it:} \emph{avia}, once more
  \item \emph{Say it:} say \emph{avia}
  \item \emph{Recall:} say the Latin for a daughter, then the Latin for a grandfather, then say \emph{avia} again
\end{itemize}

\begin{tcolorbox}[breakable,colback=teal!4,colframe=teal!35!black,title={Before you move on}]
What does avia, aviae mean? (\textquotedblleft{}A grandmother\textquotedblright{}.) Say it once more.
\end{tcolorbox}
\section[familia, familiae]{familia, familiae — a family, a household}

\label{lesson:LA-C70-familia}



\begin{quote}
Before the new one: say the Latin for a grandfather, then the Latin for a grandmother.
\end{quote}

\subsection*{You'll want to know: familia, familiae}
\textbf{familia, familiae} — \textquotedblleft{}a family, a household\textquotedblright{}.

\begin{cousinweb}
\textbf{familia} is built on \emph{famulus}, \textquotedblleft{}a household slave, a servant\textquotedblright{} (Lewis \& Short): a \emph{familia} was first everyone under one master's roof, the slaves included. English \emph{family} comes from it.
\end{cousinweb}

\begin{grammarlens}[title={What you've built}]
One more word for this chapter.
\end{grammarlens}

\subsection*{Your turn}
\begin{itemize}
\raggedright
  \item \emph{Say it:} \emph{familia}
  \item \emph{Say it:} \emph{familia}, once more
  \item \emph{Say it:} say \emph{familia}
  \item \emph{Recall:} say the Latin for a grandfather, then the Latin for a grandmother, then say \emph{familia} again
\end{itemize}

\begin{tcolorbox}[breakable,colback=teal!4,colframe=teal!35!black,title={Before you move on}]
What does familia, familiae mean? (\textquotedblleft{}A family, a household\textquotedblright{}.) Say it once more.
\end{tcolorbox}
